\section*{第 \ref{ch:g12:seq} 章 --- 数列}

\begin{solution}{exo:g12:seq:1}
设 $P(n)$ 为命题 $\sum_{k=1}^{n} k^2 = \frac{n(n+1)(2n+1)}{6}$。
\emph{基础情形：}当 $n = 0$ 时，两边都等于 $0$（空和）。
\emph{归纳步骤：}假设 $P(n)$ 成立，则
\[
\sum_{k=1}^{n+1} k^2 = \frac{n(n+1)(2n+1)}{6} + (n+1)^2
= \frac{(n+1)\bigl(n(2n+1) + 6(n+1)\bigr)}{6}
= \frac{(n+1)(2n^2 + 7n + 6)}{6}.
\]
由 $2n^2 + 7n + 6 = (n+2)(2n+3)$，上式即为
$\frac{(n+1)(n+2)(2(n+1)+1)}{6}$，即 $P(n+1)$。由归纳法，$P(n)$
对所有 $n$ 成立。
\end{solution}

\begin{solution}{exo:g12:seq:2}
$a_{n+1} - a_n = \frac{n+2}{n+1} - \frac{n+1}{n}
= \frac{n(n+2) - (n+1)^2}{n(n+1)} = \frac{-1}{n(n+1)} < 0$，故 $(a_n)$
严格递减。

$(b_n)$ 各项为正，且
$\frac{b_{n+1}}{b_n} = \frac{2^{n+1}}{n+1}\cdot\frac{n}{2^n} =
\frac{2n}{n+1} \geq 1 \iff 2n \geq n+1 \iff n \geq 1$，故 $(b_n)$
递增（$n \geq 2$ 时严格）。

$c_{n+1} - c_n = (n+1)^2 - 10(n+1) - n^2 + 10n = 2n - 9$，
当 $n \leq 4$ 时为负，当 $n \geq 5$ 时为正：$(c_n)$ 先递减
至 $c_5 = -25$（最小值），此后递增。它不单调。
\end{solution}

\begin{solution}{exo:g12:seq:3}
提出主导项：
\[
u_n = \frac{n^2(3 - 1/n + 1/n^2)}{n^2(2 + 5/n^2)} \xrightarrow[n\to+\infty]{} \frac{3}{2}.
\]
乘以共轭根式：
\[
v_n = \frac{(n+1) - n}{\sqrt{n+1} + \sqrt{n}} = \frac{1}{\sqrt{n+1}+\sqrt{n}}
\xrightarrow[n\to+\infty]{} 0.
\]
分子分母同除以 $3^n$：
\[
w_n = \frac{(2/3)^n - 1}{1 + (1/3)^n} \xrightarrow[n\to+\infty]{} \frac{0-1}{1+0} = -1,
\]
利用 $\lim q^n = 0$（当 $\abs{q} < 1$ 时）。
\end{solution}

\begin{solution}{exo:g12:seq:4}
$u_n = 5 + 3n \to +\infty$，且 $v_n = 8 \cdot (1/2)^n = 2^{3-n} \to 0$。
其部分和为
\[
\sum_{k=0}^{n} u_k = (n+1)\,\frac{5 + (5+3n)}{2} = \frac{(n+1)(10+3n)}{2}
\xrightarrow[n\to+\infty]{} +\infty,
\]
\[
\sum_{k=0}^{n} v_k = 8\,\frac{1 - (1/2)^{n+1}}{1 - 1/2}
= 16\left(1 - \left(\tfrac{1}{2}\right)^{n+1}\right)
\xrightarrow[n\to+\infty]{} 16 .
\]
\end{solution}

\begin{solution}{exo:g12:seq:5}
由 $-1 \leq \cos n \leq 1$，
\[
\frac{n-1}{n+1} \leq \frac{n + \cos n}{n+1} \leq 1,
\]
且 $\frac{n-1}{n+1} \to 1$，由夹逼定理，极限为 $1$。

对第二个极限，写作
\[
0 \leq \frac{n!}{n^n}
= \frac{1}{n}\cdot\frac{2}{n}\cdots\frac{n}{n}
\leq \frac{1}{n},
\]
因为每个因子 $\frac{k}{n}$（$2 \leq k \leq n$）都不超过 $1$。
由 $\frac1n \to 0$，夹逼定理给出 $\frac{n!}{n^n} \to 0$。（
题目提示的几何估计同样可行：当 $k \leq n/2$ 时，每个因子至多为
$\frac12$，从而得到更强的界 $(1/2)^{\floor{n/2}}$。）
\end{solution}

\begin{solution}{exo:g12:seq:6}
\emph{1.} $u_0 = 0 \in \intcc{0}{2}$。若 $0 \leq u_n \leq 2$，则
$2 \leq u_n + 2 \leq 4$，故 $\sqrt{2} \leq u_{n+1} \leq 2$；特别地
$0 \leq u_{n+1} \leq 2$。由归纳法，该性质对所有 $n$ 成立。

\emph{2.} $u_{n+1} - u_n = \sqrt{u_n + 2} - u_n$。对 $x \in \intcc{0}{2}$，
$\sqrt{x+2} \geq x \iff x + 2 \geq x^2 \iff (2-x)(x+1) \geq 0$，显然成立。
因此 $(u_n)$ 递增。

\emph{3.} 递增且以 $2$ 为上界的数列 $(u_n)$ 收敛于某个
$\ell \in \intcc{0}{2}$。在 $u_{n+1} = \sqrt{u_n + 2}$ 中取极限
（映射 $x \mapsto \sqrt{x+2}$ 连续），得 $\ell = \sqrt{\ell + 2}$，
故 $\ell^2 - \ell - 2 = 0$，即 $\ell \in \{-1, 2\}$。由 $\ell \geq 0$，所以
$\lim u_n = 2$.
\end{solution}

\begin{solution}{exo:g12:seq:7}
\emph{1.} 两次服药之间有 $40\%$ 的药物被排出，因此
药量 $u_n$ 变为 $0.6\,u_n$；下一次服药增加 $1$ 个单位：
$u_{n+1} = 0.6\,u_n + 1$.

\emph{2.} $v_{n+1} = u_{n+1} - 2.5 = 0.6\,u_n + 1 - 2.5 = 0.6(u_n - 2.5)
= 0.6\,v_n$：$(v_n)$ 为公比 $0.6$ 的等比数列，首项
为 $v_0 = 1 - 2.5 = -1.5$。由此 $v_n = -1.5 \times 0.6^n$，且
\[
u_n = 2.5 - 1.5 \times 0.6^n .
\]

\emph{3.} 由 $0.6^n \to 0$，得 $u_n \to 2.5$：药物量稳定在
$2.5$ 个单位。
\end{solution}

\begin{solution}{exo:g12:seq:8}
\emph{1.} $u_0 = 3 > 1$。若 $u_n > 1$，则 $u_n + 2 > 0$，且
\[
u_{n+1} - 1 = \frac{4u_n - 1 - u_n - 2}{u_n + 2} = \frac{3(u_n - 1)}{u_n + 2} > 0 .
\]
由归纳法，$u_n > 1$ 对所有 $n$ 成立（特别地 $u_n + 2 \neq 0$，
故该数列有定义）。

\emph{2.} 利用上式，
\[
v_{n+1} = \frac{1}{u_{n+1} - 1} = \frac{u_n + 2}{3(u_n - 1)}
= \frac{(u_n - 1) + 3}{3(u_n - 1)} = \frac{1}{3} + v_n .
\]
故 $(v_n)$ 为公差为 $\frac13$ 的等差数列，且
首项 $v_0 = \frac{1}{u_0 - 1} = \frac12$。

\emph{3.} 由 $v_n = \frac12 + \frac{n}{3}$，得
$u_n = 1 + \frac{1}{v_n} = 1 + \frac{6}{3 + 2n}$。由 $v_n \to +\infty$，
故 $u_n \to 1$。
\end{solution}

\begin{solution}{exo:g12:seq:9}
\emph{1.} $H_{2n} - H_n = \sum_{k=n+1}^{2n} \frac{1}{k}$ 为 $n$ 项之和，
每项均不小于 $\frac{1}{2n}$，因此
$H_{2n} - H_n \geq n \cdot \frac{1}{2n} = \frac12$。

\emph{2.} 对 $k$ 用归纳法：$H_{2^0} = H_1 = 1 \geq 1$。若
$H_{2^k} \geq 1 + \frac{k}{2}$，则由 1 的结论取 $n = 2^k$，
\[
H_{2^{k+1}} \geq H_{2^k} + \frac12 \geq 1 + \frac{k+1}{2}.
\]
数列 $(H_n)$ 递增（每一步增加 $\frac{1}{n+1} > 0$），且
子列 $H_{2^k}$ 无界，故 $(H_n)$ 无上界。
递增且无界，故其趋于 $+\infty$
（见\cref{thm:g12:seq:monotone}）。
\end{solution}

\begin{solution}{exo:g12:seq:10}
\emph{1.} 数列 $d_n = b_n - a_n$ 满足
$d_{n+1} - d_n = (b_{n+1} - b_n) - (a_{n+1} - a_n) \leq 0$，故 $(d_n)$
递减；由 $d_n \to 0$，得 $d_n \geq 0$ 对所有 $n$ 成立（若递减数列出现负项，
则其此后各项恒位于该负项之下，从而无法
趋于 $0$）。因此 $a_n \leq b_n$。

\emph{2.} 由 $a_n \leq b_n \leq b_0$，递增数列 $(a_n)$
有上界，故收敛于某数 $\ell$。类似地，递减数列 $(b_n)$
有下界 $a_0$，收敛于某数 $\ell'$。于是
$\ell' - \ell = \lim (b_n - a_n) = 0$，故 $\ell = \ell'$。

\emph{3.} $(a_n)$ 为（严格）递增数列，因为
$a_{n+1} - a_n = \frac{1}{(n+1)!} > 0$。对 $(b_n)$，
\[
b_{n+1} - b_n = \frac{1}{(n+1)!} + \frac{1}{(n+1)(n+1)!} - \frac{1}{n\,n!}
= \frac{n(n+1) + n - (n+1)^2}{n(n+1)(n+1)!}
= \frac{-1}{n(n+1)(n+1)!} < 0 ,
\]
故 $(b_n)$ 递减。最后，$b_n - a_n = \frac{1}{n\,n!} \to 0$。两
数列相邻，因此收敛于同一极限。
\end{solution}

\begin{solution}{pb:g12:seq:1}
\textbf{1.} 当 $n = 1$ 时成立（$1 = 1^2$）。若
$1 + 3 + \dots + (2n - 1) = n^2$，则加上下一个奇数
$n^2 + (2n + 1) = (n + 1)^2$，即递推成立。由归纳法，
对所有 $n \geq 1$ 成立。

\textbf{2.} $2^0 = 1 > 0$。若 $2^n > n$，则
$2^{n+1} = 2 \cdot 2^n > 2n \geq n + 1$（对 $n \geq 1$ 成立，
$n = 0$ 时可直接验证）。递推成立，结论得证。

\textbf{3.} $n = 0$ 时，$1 \geq 1$。若
$(1 + x)^n \geq 1 + nx$，两边乘以 $1 + x \geq 1 > 0$：
$(1 + x)^{n+1} \geq (1 + nx)(1 + x) = 1 + (n + 1)x + nx^2
\geq 1 + (n + 1)x$。

\textbf{4.} 从 $n = 1$ 到 $n = 2$ 这一步：在两颗弹珠中，
“前 $n$ 颗”与“后 $n$ 颗”分别只是两颗
\emph{互不相交}的单颗弹珠 --- 没有公共弹珠连接这两组，
因而无法保证它们的颜色一致。递推论证
隐含地要求两组有重叠，
这仅在 $n \geq 2$ 时成立；取基础情形 $n = 1$ 时，
推理链条无法开始。

\textbf{5.} $4^0 - 1 = 0 = 3 \times 0$。若 $4^n - 1 = 3k$，
则 $4^{n+1} - 1 = 4(4^n - 1) + 3 = 3(4k + 1)$，递推成立。

\textbf{6.} $x_1 = \frac32$，$x_2 = \frac{17}{12}$，
$x_3 = \frac{577}{408}$。

\textbf{7.} $x_{n+1}^2 - 2 = \frac{(x_n^2 + 2)^2 - 8x_n^2}
{4x_n^2} = \frac{(x_n^2 - 2)^2}{4 x_n^2}$：平方式除以一个
正数，故 $\geq 0$，且 $> 0$ 当 $x_n^2 \neq 2$ 时成立。
归纳法：$x_0 = 2 > 0$ 且 $x_0^2 = 4 > 2$；若 $x_n > 0$
且 $x_n^2 > 2$，则 $x_{n+1}$（正数的平均）为
正，且 $x_{n+1}^2 - 2 > 0$。

\textbf{8.} $x_{n+1} - x_n = \frac{2 - x_n^2}{2x_n} < 0$，由
第 7 题，数列严格递减。

\textbf{9.} 数列递减且有下界（可取下界 $1$，因为
$x_n^2 > 2 > 1$ 且 $x_n > 0$）。由单调收敛
定理，$(x_n)$ 收敛于某数 $L \geq 1$。

\textbf{10.} 极限与代数运算相容：由
$x_{n+1} = \frac12\left(x_n + \frac{2}{x_n}\right)$ 及
$x_n \to L \geq 1 > 0$：
得 $L = \frac12\left(L + \frac2L\right)$，故 $L^2 = 2$，又 $L$
为正，故 $L = \sqrt2$。结论：收敛性已证，
极限已确定 --- 海伦以优异成绩无罪开释。

\textbf{11.} $x_{n+1} - \sqrt2 = \frac{x_n^2 - 2\sqrt2\,x_n +
2}{2x_n} = \frac{(x_n - \sqrt2)^2}{2x_n}$：即
$e_{n+1} = \frac{e_n^2}{2x_n}$，又由 $x_n > \sqrt2$ 可得
$e_{n+1} \leq \frac{e_n^2}{2\sqrt2}$。平方误差：每一步
使正确的小数位数翻倍，这一点自
九年级起便已被观察到。

\textbf{12.} $e_0 \approx 0.5858$, $e_1 \approx 0.0858$,
$e_2 \approx 0.00245$, $e_3 \approx 2.1 \times 10^{-6}$.
比值 $\frac{e_1}{e_0^2} \approx 0.25 = \frac{1}{2x_0}$；
$\frac{e_2}{e_1^2} \approx 0.333 = \frac{1}{2x_1}$;
$\frac{e_3}{e_2^2} \approx 0.353 \approx \frac{1}{2x_2}$：这正是
定理的实际应用。

\textbf{13.} $H_{2^{18}} \geq 1 + 9 = 10$：约
$260\,000$ 项（$2^{18} = 262\,144$），仅为超过 $10$ ---
发散慢如爬行（而 $H_n > 100$ 所需的项数
将多于任何一座图书馆中的原子数）。

\textbf{14.} $S_n = 2 - \frac{1}{2^n}$（等比数列求和），且
$\frac{1}{2^n} \to 0$（见\cref{thm:g12:seq:geometric}），故
$S_n \to 2$：被无尽咬去的巧克力棒趋于
却永远达不到整块 --- 如今以
极限的正式语言来表述。

\textbf{15.} 对 $k \geq 2$：
$\frac{1}{k^2} \leq \frac{1}{k(k-1)} = \frac{1}{k-1} -
\frac1k$，于是
$S_n \leq 1 + \left(1 - \frac1n\right) < 2$：递增且
有上界，从而收敛（单调收敛定理）。欧拉
后来算出该极限为 $\frac{\pi^2}{6}$。

\textbf{16.} 项趋于 $0$ 是和式收敛的 \emph{必要}条件，
但并不能决定和式的敛散性：调和数列的项
$\frac1n \to 0$，和式却发散；
$\frac{1}{n^2} \to 0$，和式却收敛。项趋于零的 \emph{快慢}
才是问题的全部 --- 级数理论
将在大学分册中建立。

\textbf{17.} $a_1 = \sqrt2 \approx 1.41421$, $b_1 = 1.5$;
$a_2 \approx 1.45648$，$b_2 \approx 1.45711$：两次迭代
已在小数点后三位上一致 --- 速度惊人。

\textbf{18.} $b_{n+1} - a_{n+1} = \frac{a_n + b_n}{2} -
\sqrt{a_n b_n} = \frac{(\sqrt{b_n} - \sqrt{a_n})^2}{2} \geq
0$：两个平均数保持大小有序。$(a_n)$ 递增：
$a_{n+1} = \sqrt{a_n b_n} \geq \sqrt{a_n \cdot a_n} = a_n$;
$(b_n)$ 对称地递减。

\textbf{19.} $\frac{b_{n+1} - a_{n+1}}{b_n - a_n} =
\frac{(\sqrt{b_n} - \sqrt{a_n})^2}
{2(\sqrt{b_n} - \sqrt{a_n})(\sqrt{b_n} + \sqrt{a_n})}
= \frac{\sqrt{b_n} - \sqrt{a_n}}{2(\sqrt{b_n} + \sqrt{a_n})}
\leq \frac12$：间距至少减半，故 $b_n - a_n \to 0$；
结合第 18 题，两数列相邻，其共同极限为
$M(1, 2)$。

\textbf{20.} 第三次迭代给出
$a_3 \approx b_3 \approx 1.456791$：$M(1, 2) \approx
1.456791$，仅需三次摇柄即可（间距大致
以 \emph{平方} 速度缩小，与海伦法相仿）。本题诸数列的速度排序，
由慢到快依次为：调和数列的部分和（冰川般
的缓慢发散），等比数列的部分和（误差每步减半），
海伦法与算术几何平均（误差每步平方） --- 正是
算术几何平均的超凡速度，使高斯意识到
自己触及了分析学的一条新矿脉。
\end{solution}
